% !PRINTED 609
\weberpage{609}
Now if
\weq{15}{\alpha_1 = a_{1,1}, \quad \alpha_2 = a_{1,2}+a_{2,2}\Theta}
is a basis of $\varphi$, then on the other hand (by \S.\,146)
\weq{16}{N_a(\varphi) = a_{1,1}a_{2,2},}
and $a_{1,1}$ is the least positive integral rational number that is divisible by $\varphi$.

Now we can show, on the other hand, that this number is $= ec$, from which $ec = a_{1,1}$ follows, and, by (14) and (16), $e = a_{2,2}$.

For if an integral rational number $m$ is to be divisible by $\varphi$, then, if $\omega'$ is conjugate to $\omega$, since
\[
ct_1^2 + bt_1t_2 - at_2^2
\]
is a unit functional,
\[
\frac{m(ct_1^2+bt_1t_2-at_2^2)}{\varphi} = \frac{m}{e}(t_1+\omega' t_2)
\]
must be an integral functional. Consequently
\[
\frac{m}{e} \quad \text{and} \quad \frac{m}{e}\omega'
\]
must be integers, from which it follows that $m : e$ is an integral rational number divisible by $c$, and hence $m$ one divisible by $ec$. Since, on the other hand, $ec$ is divisible by $\varphi$, it follows that $ec$ is the least number satisfying these requirements. Moreover, it also results that the integer $ec\omega$ is divisible by $\varphi$. For in the representation (12), $ec\omega$ is a coefficient of the integral function $\varphi$, and is therefore divisible by $\varphi$ (\S.\,142).

If we express $\Theta$ in (15) in terms of $\omega$, then by (11) we obtain
\[
\begin{aligned}
\alpha_2 &= a_{1,2} - \frac{eb}{2} + ec\omega, \quad b \text{ even}\\
&= a_{1,2} - \frac{e(b-1)}{2} + ec\omega, \quad b \text{ odd,}
\end{aligned}
\]
and thus
\[
a_{1,2} \equiv \frac{eb}{2} \quad \text{or} \quad \equiv \frac{e(b-1)}{2}
\]
with respect to the modulus $\varphi$, and consequently, since both sides are rational, also with respect to the modulus $a_{1,1}=ec$. Accordingly $a_{1,2}$ may be set equal to the value $\frac{1}{2}eb$ or $\frac{1}{2}e(b-1)$, and we obtain $\alpha_2 = ec\omega$, which was to be proved.

In order to give a clear summary of what has hereby been proved, we may state the following theorems:
%CONT>>
